% Paper 5, §12 — Watermarking and provenance
% Anchors: kb/notes/alignment/watermarking-and-provenance.md
%          kb/excerpts/kirchenbauer2023-watermark.md

\section{Watermarking and provenance}
\label{sec:watermarking}

The threats covered in \cref{sec:scheming,sec:alignment-faking} concern
what the model \emph{does}; the threat in this section concerns what
\emph{has been done} --- the question of whether a given span of text
was produced by a particular language model at all. Watermarking is the
sampling-time intervention that embeds a statistical signal into the
token stream, designed to be invisible to a casual human reader and
algorithmically detectable from a short keyed test. Provenance is the
broader project of certifying the origin of generated content;
watermarking is the technical primitive most-discussed for text. The
thesis of this section is that the canonical Kirchenbauer 2023
construction --- a hash-keyed green/red partition of the vocabulary at
each step, with a soft logit nudge toward the green list, scored by a
per-token z-statistic --- has the \emph{cryptographic} part of the
problem solved. The 2024 production deployments (Google's
SynthID-Text, OpenAI's deferred deployment) carry that construction
forward. What the field does not yet have is a watermark robust to a
determined removal attempt: paraphrase, fine-tuning, and collusion
each define an attack surface where the signal degrades by amounts
that vendors and academics characterize differently.

\subsection{The green/red-list construction}
\label{sec:watermarking-construction}

\citet{kirchenbauer2023-watermark} fix the canonical construction%
\footnote{KB excerpt: \texttt{kb/excerpts/kirchenbauer2023-watermark.md\#sec-greenred}.}
\citep[\S{}abstract]{kirchenbauer2023-watermark}:

\begin{quote}\itshape
The watermark works by selecting a randomized set of ``green'' tokens
before a word is generated, and then softly promoting use of green
tokens during sampling.
\end{quote}

The construction is stateless beyond a small hash window. Let $V$ be
the vocabulary, $\gamma \in (0, 1)$ a partition fraction, and
$\delta > 0$ a logit-bias constant. At generation step $t$, with
preceding tokens $s^{(1)}, \ldots, s^{(t-1)}$:

\begin{enumerate}
\item Hash a fixed-width suffix --- in the simplest variant
$s^{(t-1)}$ alone --- through a keyed pseudo-random function to seed
an RNG.
\item Partition $V$ into a green list $G_t \subset V$ of size
$\gamma|V|$ and a red list $V \setminus G_t$ of size
$(1 - \gamma)|V|$.
\item Add $\delta$ to the logit of every $v \in G_t$ before the
softmax: $\tilde{\ell}_v = \ell_v + \delta \cdot \mathbb{1}[v \in G_t]$.
\item Sample $s^{(t)}$ from the modified distribution.
\end{enumerate}

The default settings reported in the paper are $\gamma = 0.25$ and
$\delta = 2.0$%
\footnote{KB excerpt: \texttt{kb/excerpts/kirchenbauer2023-watermark.md\#sec-greenred}.}.
Three properties of this construction are load-bearing for everything
that follows. First, the partition $G_t$ depends only on the keyed hash
of a small recent window, so generation has $O(|V|)$ overhead per step
and detection is parallelizable. Second, $\delta$ is added before the
softmax, so the scheme is a \emph{soft} promotion: low-entropy positions
where one token's logit dominates by more than $\delta$ are effectively
unaffected, and high-entropy positions get a measurable nudge. Third,
the green list is determined by the secret key plus the candidate text,
not by any property of the model --- detection needs the keyed hash
function and the candidate string, nothing more
\citep[\S{}abstract]{kirchenbauer2023-watermark}:

\begin{quote}\itshape
The watermark can be embedded with negligible impact on text quality,
and can be detected using an efficient open-source algorithm without
access to the language model API or parameters.
\end{quote}

This third property --- \emph{detector-only-needs-the-key} --- is what
distinguishes the Kirchenbauer scheme from earlier classifier-based
AI-text-detection efforts. The detector is a fixed statistical test,
not another model. False-positive rates are derived from the test
statistic, not measured on a held-out set.

\subsection{Detection: the per-token z-score}
\label{sec:watermarking-detection}

Given a candidate text $s$ of length $T$, the detector recomputes the
green list at each position from the keyed hash and counts the number
of green-listed tokens $|s|_G$. Under the null hypothesis that $s$ was
not watermarked, each token is in $G_t$ independently with probability
$\gamma$, so $|s|_G$ is binomially distributed with mean $\gamma T$ and
variance $T \gamma (1 - \gamma)$. The normalized test statistic is the
one-sided z-score%
\footnote{KB excerpt: \texttt{kb/excerpts/kirchenbauer2023-watermark.md\#sec-detection}.}
\citep[\S{}sec-detection]{kirchenbauer2023-watermark}:
\[
z \;=\; \frac{|s|_G - \gamma T}{\sqrt{T \gamma (1 - \gamma)}}.
\]
A large $z$ rejects the null with a one-sided $p$-value computed
directly from the normal CDF. At $\gamma = 0.25$, a span of $T = 200$
tokens with $|s|_G = 100$ greens (against a null mean of $50$) yields
$z = (100 - 50)/\sqrt{200 \cdot 0.25 \cdot 0.75} \approx 8.16$,
corresponding to a one-sided $p < 10^{-15}$ --- comfortably below any
deployment-relevant false-positive budget. The paper exposes the
$p$-value as the deployment knob: a publisher willing to tolerate one
false positive in $10^6$ pieces of long-form content sets the
$z$-threshold accordingly and reads off the minimum span required for
detection at a given watermark strength.

The same calculation gives the false-positive rate (FPR) directly: at
threshold $z^\star$, $\mathrm{FPR} = 1 - \Phi(z^\star)$, where $\Phi$
is the standard normal CDF. This is the structurally cleanest piece of
the construction: FPR control is by the test statistic alone, not by
held-out data, so a deployer can guarantee FPR without access to the
generation distribution at all.

\subsection{The spike-entropy detection bound}
\label{sec:watermarking-spike-entropy}

The detection-rate analysis is where the construction earns its
information-theoretic frame. The paper%
\footnote{KB excerpt: \texttt{kb/excerpts/kirchenbauer2023-watermark.md\#abstract};
the precise constants in the bound are pending direct PDF anchor
verification.}
\citep[\S{}abstract]{kirchenbauer2023-watermark}:

\begin{quote}\itshape
We propose a statistical test for detecting the watermark with
interpretable p-values, and derive an information-theoretic framework
for analyzing the sensitivity of the watermark.
\end{quote}

The sensitivity result that organizes \cref{sec:watermarking-construction}
is that detection power scales not with the Shannon entropy of the
model's distribution, but with a sharper functional --- the
\emph{spike entropy} at parameter $\gamma$. For a distribution
$p$ over $V$, define%
\footnote{The functional form is the one cited in
\texttt{kb/notes/alignment/watermarking-and-provenance.md} \S 1; the
exact paper-side anchor is flagged for the deepening pass.}
\deepencite{kirchenbauer2023-watermark §sensitivity-bound}:
\[
S_\gamma(p) \;=\; \sum_{v \in V} \frac{p_v}{1 + \gamma' \, p_v},
\]
with $\gamma'$ a constant tied to the partition fraction $\gamma$ and
the bias $\delta$ in the manner the sensitivity analysis specifies. The
substantive content of the bound, sketched here at the form level
rather than at the constant level: the expected fraction of green
tokens in the watermarked sample is
\[
\E[\,|s|_G \,/\, T\,] \;=\; \gamma + c(\gamma, \delta) \cdot S_\gamma(p) \cdot (1 - \gamma) + o(\cdot),
\]
with $c(\gamma, \delta)$ an increasing function of $\delta$ that the
paper computes. The qualitative claim is that detection accumulates
faster on samples drawn from \emph{spiky} distributions (low spike
entropy, one or a few tokens carrying most of the mass) than on
samples drawn from flat distributions, and that this dependence is
captured by $S_\gamma$ rather than by the standard Shannon
$H(p) = -\sum_v p_v \log p_v$. The two functionals can disagree: a
distribution with high Shannon entropy that is nevertheless
concentrated on a moderate number of green-list tokens watermarks
better than the Shannon entropy alone would predict.

\begin{intuition}
The watermark is a stochastic side-channel that the model itself
supplies. At each step, the model's softmax expresses how much
flexibility it has --- how many distinct tokens are roughly
acceptable in this position. Watermarking taxes that flexibility by a
constant $\delta$ in favor of a hash-keyed green list. At a low-entropy
position --- a factual lookup, a code keyword, a closing punctuation
mark --- there is no flexibility for the watermark to spend, and the
nudge does nothing. At a high-entropy position --- a stylistic
adverb, an open-ended sentence continuation --- the watermark is
nearly free, because the model would have been almost indifferent
between the green and red alternatives. The detection bound is the
quantitative statement of this picture: the watermark accumulates
exactly where the entropy is, and the right entropy functional is
$S_\gamma$, the entropy reweighted by how much logit-bias $\delta$
can move per token. (Returning to canonical form: the per-token
z-score in \cref{sec:watermarking-detection} is the statistic; the
$S_\gamma$-dependence is the rate at which it accumulates over $T$.)
\end{intuition}

The implication for deployment is that watermarking is not uniformly
applicable across content types. Long-form prose, creative writing,
and conversational responses carry high spike entropy in most
positions and are watermarkable at the per-paragraph scale. Code,
verbatim citation, structured data extraction, and short factual
answers carry low spike entropy and are effectively unwatermarkable in
their generation regions --- the green-list nudge cannot be
distinguished from the model's natural sampling%
\footnote{KB note: \texttt{kb/notes/alignment/watermarking-and-provenance.md\#1}
\S~1 (\textsc{intuition}).}.

\subsection{Soft promotion and the distortion-detectability tradeoff}
\label{sec:watermarking-soft}

The choice of a $\delta$-additive logit bias rather than a hard
partition is the design decision that makes the rest workable. The
paper%
\footnote{KB excerpt: \texttt{kb/excerpts/kirchenbauer2023-watermark.md\#sec-soft}.}
discusses three failure modes of a hard partition that restricts
generation to $G_t$ alone: it heavily distorts high-entropy
generation by forcing the model away from its natural distribution,
it makes low-entropy positions impossible to watermark because no
green-list-friendly alternative exists, and it is trivially detectable
by frequency analysis even without the key. The soft scheme avoids
all three by accepting a quantitative tradeoff: larger $\delta$ gives
a stronger watermark per token but introduces a measurable bias in
the marginal output distribution. The paper reports
``negligible impact on text quality'' at the default
$(\gamma, \delta) = (0.25, 2.0)$
\citep[\S{}abstract]{kirchenbauer2023-watermark}; deployment is a
choice of where to sit on this curve.

The same soft-promotion intuition explains why hard-partition
watermarks have not survived as the field-default. A scheme that
forces green-only generation gives the strongest possible per-token
signal but breaks at exactly the positions deployers care about ---
factual answers, code, citations, structured outputs --- because at
those positions the green list either contains the right answer (in
which case the scheme reduces to passthrough) or does not (in which
case the scheme either refuses to generate or generates a wrong
answer). The $\delta$-additive form lets the scheme degrade
gracefully at low-entropy positions and accumulate signal at
high-entropy ones, which is the empirical operating point modern
deployments target.

\subsection{Production: SynthID-Text and OpenAI's deferred deployment}
\label{sec:watermarking-production}

By 2024, two production-scale text-watermarking efforts had landed in
public conversation. Google DeepMind's \textbf{SynthID-Text} (Dathathri
et al., \emph{Nature} 2024) replaces the logit-bias step with a
\emph{tournament} over candidate tokens%
\footnote{KB note: \texttt{kb/notes/alignment/watermarking-and-provenance.md\#3-synthid-text-dathathri-et-al-2024-nature}.}
\deepencite{dathathri2024-synthid §method, citation-pending}: at each
step the sampler draws $N$ candidate next-tokens, pairs them, evaluates
a hash-derived watermark function $g$ on each, and advances the higher
$g$-value candidate; after $\log_2 N$ rounds, a single token wins. The
detection statistic averages $g$-values across the generated tokens and
flags spans whose average is above the no-watermark baseline. Two modes
are reported: a \emph{distortionary} mode that visibly biases the
output distribution and a \emph{non-distortionary} mode that, under
a sufficient-candidate-diversity assumption, leaves the marginal
distribution statistically unchanged \deepencite{dathathri2024-synthid
§distortion-modes, citation-pending}. SynthID-Text is reported as
deployed in the Gemini App, the first production-scale text-watermark
deployment by a frontier lab; detection code accompanied the
\emph{Nature} publication%
\footnote{KB note: \texttt{kb/notes/alignment/watermarking-and-provenance.md\#3-3-production-deployment}.}.

OpenAI publicly described a watermark prototype in 2023--2024 but
\emph{deferred deployment} on the basis of robustness and policy
concerns; the available account is a vendor blog post rather than a
peer-reviewed publication, so this paper treats it as a tier-B signal
rather than a tier-A claim
\deepencite{openai2024-watermark-deferral, citation-pending}.
The 2024 picture is therefore one production system shipping a
Kirchenbauer-lineage watermark at scale, one large vendor with a
working prototype not deployed, and a several-arXiv-paper academic
literature on the spike-entropy / paraphrase / fine-tune attack
surface. The Kirchenbauer construction is the common ancestor of the
shipped scheme and the deferred one.

\subsection{The attack surface}
\label{sec:watermarking-attacks}

Three named attacks define the open frontier: paraphrase, fine-tuning,
and collusion. Each pressures a different invariant of the construction.

\textbf{Paraphrase attacks} round-trip the watermarked text through a
different model (or a non-watermarked instance of the same model) and
re-emit it. Kirchenbauer's own follow-up
\deepencite{kirchenbauer2024-reliability §paraphrase, citation-pending}
reports that even after strong human paraphrasing, the watermark
remains detectable at $T \approx 800$ tokens at a $10^{-5}$
false-positive rate, because paraphrasing preserves enough short
n-gram structure that the green-list hash continues to land on green
tokens at above-chance rates --- just slowly. The countervailing
position from \citet{sadasivan2023-detection}
\deepencite{sadasivan2023-detection §main-result, citation-pending}
argues that any watermark detectable from \emph{short} spans must be
removable by a paraphrase attacker with a comparable-quality model;
the disagreement is over what ``deployable'' requires, not over the
asymptotic detection rate%
\footnote{KB note: \texttt{kb/notes/alignment/watermarking-and-provenance.md\#2-2-the-fundamental-limit}
\S~2.2.}.

\textbf{Fine-tuning attacks} retrain the watermarked model (or a base
model) on a small corpus of unrelated text, with the goal of disrupting
whichever internal feature the watermark relies on. For a green/red
construction, the watermark is not a model feature --- it is a
sampling-time intervention --- so fine-tuning the base weights does
not directly remove the keyed hash. What fine-tuning can disrupt is
the \emph{deployment context} that applies the intervention: a
fine-tuned model derived from a watermarked model's checkpoints loses
the watermark unless the fine-tuned deployer also runs the
green-list sampler. The vendor and academic claims diverge on what
fine-tune budget defeats which schemes
\deepencite{watermark-fine-tune-survey, citation-pending}, with the
disagreement structured around whether the relevant attack surface is
the keyed sampling layer (in which case fine-tuning is irrelevant) or
the deployer's pipeline (in which case fine-tuning to a fork that
omits the sampler breaks the watermark trivially).

\textbf{Collusion attacks} mix watermarked outputs with unwatermarked
text from another source, diluting the green-token rate below the
detection threshold. The detector's z-score is the per-token statistic
of \cref{sec:watermarking-detection}: at $\gamma = 0.25$, a $50$/$50$
mix of watermarked and unwatermarked text has expected green rate
$0.5 \cdot E_{\mathrm{wm}} + 0.5 \cdot \gamma$, where $E_{\mathrm{wm}}$
is the watermark's natural green rate. The detector applied to the
full mixed span sees a degraded $z$ that scales with the mixed-span
length and the mixing fraction; if the deployer can segment the text
into pure-watermarked spans, detection recovers, but the segmentation
problem is itself adversarial%
\footnote{KB note: \texttt{kb/notes/alignment/watermarking-and-provenance.md\#5}
\S~5 (frontier).}.

\begin{contradiction}
\textbf{Watermarking robustness against fine-tuning.} The Kirchenbauer
2024 reliability paper \deepencite{kirchenbauer2024-reliability
§paraphrase-and-finetune, citation-pending} reports that the
green-list watermark survives moderate fine-tuning at the cost of a
slower detection accumulation rate, and the Google SynthID-Text
deployment account
\deepencite{dathathri2024-synthid §robustness, citation-pending}
makes a similar claim for the tournament-sampling variant. Independent
adversarial analyses --- notably the SRI Lab ETH probing study of
SynthID%
\footnote{KB note: \texttt{kb/notes/alignment/watermarking-and-provenance.md\#3-3-production-deployment}
\S~3.3 (\textsc{contradiction}).}
\deepencite{eth-synthid-probe-2024, citation-pending}
--- find that aggressive paraphrase and fine-tune budgets do defeat
deployed schemes in the short-form regime. The two camps are not
numerically far apart on the asymptotic rate. They disagree on whether
the threat model that matters is a casual user trying to bypass an
enterprise content policy (vendor framing, where the watermark holds)
or a determined adversary with a comparable-quality model and a
modest fine-tune budget (academic framing, where the watermark
degrades rapidly). As of 2026-Q2, the field does not yet have a
public adversarial benchmark where attacker-fine-tune-budget is the
$x$-axis and detector-AUC is the $y$-axis with multiple watermark
schemes plotted together; \cref{sec:contradictions} flags this as
the evidence that would close the dispute.
\end{contradiction}

\subsection{Limits and forward link}
\label{sec:watermarking-forward}

What watermarking can and cannot do is now empirically clear in
outline. It can supply discovery-grade provenance on long-form
prose --- enough signal that, given a few hundred tokens and a keyed
detector, a publisher can reject AI-generated content with a tunable
false-positive rate. It can plausibly support compliance-style
``marking'' obligations under regulations like the EU AI Act
Article 50 transparency requirement
\deepencite{eu-ai-act-art-50, citation-pending}, where the requirement
is that synthetic content carry a marking, not that the marking
withstand adversarial removal. What it cannot yet do is forensic
attribution under adversarial pressure on short text, robust survival
through arbitrary paraphrase by a comparable-quality model, or
cross-vendor interoperability the way image-and-audio C2PA
provenance does for non-text media. Open-weights models are an
additional structural limit: a user running an open-weights model
locally can simply disable the watermark sampler, so coverage of the
deployed watermark fleet caps at the closed-source share of generated
text \deepencite{open-weights-watermark-coverage-2024, citation-pending}
--- and that share is shrinking as open weights catch up to frontier
closed-weights at the capability axes a deployer would actually use.

The structural location of watermarking in the layered defense view
is the deployment-side adjacency: like \cref{sec:scalable-oversight}'s
debate / scalable-oversight protocols, watermarking is an intervention
applied at or after generation time, not a property of the model's
training. The two share the constraint that they have to work
\emph{against} a model whose capability is still rising. The
contradiction surfaced in this section --- the divergence between
vendor-friendly and academic-adversarial fine-tune-robustness claims
--- joins the alignment-evaluation contradictions of
\cref{sec:scheming,sec:alignment-faking} and the eval-instrument
contradictions of \cref{sec:knowledge-benchmarks,sec:adversarial-robustness}
in the concentrated discussion of \cref{sec:contradictions},
density-ranked alongside the others as one of the open questions
the next round of evaluation infrastructure has to attack directly.
